\chapter{Linear algebra}\label{ch:linear_algebra}

Linear algebra is a branch of mathematics that is both very accessible and enormously useful. It studies \hyperref[def:vector_space]{vector spaces}, mostly over \hyperref[def:real_numbers]{real} or \hyperref[def:complex_numbers]{complex} numbers, and \hyperref[def:linear_function]{linear maps} between them. For \hyperref[thm:vector_space_dimension]{finite-dimensional} vector spaces, this reduces to studying \hyperref[def:array/matrix]{matrices}, which also leads to a very rich computational theory.

From an algebraic perspective, we have discussed some basics of vector spaces in \fullref{sec:modules} and \fullref{sec:vector_spaces}. The key takeaways are:
\begin{itemize}
  \item Abstract vector spaces, as defined in \cref{def:vector_space} with all related notions listed inside --- \hyperref[def:linear_function]{linear functions}, \hyperref[def:linear_span]{linear spans} and \hyperref[def:vector_space/quotient]{quotients}, as well as the \hyperref[def:category]{category} \( \cat[\BbbK]{Vect} \) of vector spaces over \( \BbbK \).

  \item Linear combinations, as defined in \cref{def:linear_combination}, providing a notion of abstract syntax in vector spaces.

  \item Hamel bases, defined in \cref{def:hamel_basis} and \cref{def:ordered_hamel_basis}, with the corresponding decomposition defined in \cref{def:basis_decomposition} and the dimensions defined in \cref{def:vector_space_dimension}.

  \item \Fullref{thm:rank_nullity_theorem}.
\end{itemize}
